\chapter{Requirements and System Analysis}
\label{chap:requirements}

This chapter defines what the system must do and how each requirement can be tested.

\section{Baseline System and Scope}

The starting point is an existing pipeline for Python notebooks stored on GitHub. It reads repository and notebook information from SQLite, finds dependencies, creates a separate Python environment, executes notebooks, compares new and saved outputs, and stores the results. FAIR Jupyter exports this database information and uses it to build an RDF knowledge graph. This thesis extends those components instead of creating a separate system.

The process starts with a resource URL or a resource already listed in the database. It ends with database results, notebook categories, and an RDF graph that can be rebuilt. The extension must recognize GitHub, Codeberg, and Zenodo; download their content; convert their metadata into common fields; execute Python notebooks through the same main route; and record where each resource came from. It must also classify notebook purpose so the results can be compared by platform and category.

\begin{table}[H]
\centering
\caption{Scope of the thesis system.}
\label{tab:scope}
\footnotesize
\begin{tabularx}{\textwidth}{@{}XX@{}}
\toprule
Included & Outside the present scope\tabularnewline
\midrule
Public GitHub and Codeberg repositories & Private repositories requiring user-specific access\tabularnewline
Public Zenodo records with downloadable files & Arbitrary research repositories beyond Zenodo\tabularnewline
Python notebook environment reconstruction & Guaranteed reproduction of every operating system\tabularnewline
Execution and stored-output comparison & Validation of the scientific conclusion\tabularnewline
Normalized descriptive metadata & Automatic repair of missing source metadata\tabularnewline
Rule-based and AI-model classification & Treating automated labels as human ground truth\tabularnewline
RDF materialisation and SPARQL validation & Operation of a permanent public triple-store service\tabularnewline
\bottomrule
\end{tabularx}
\end{table}

The original database is kept read-only because it contains the starting research data. The pipeline copies it to \texttt{data/output/db/} and writes all new columns and results to that copy. Existing GitHub rows must still work after the database structure is extended.

A notebook that runs successfully is not automatically scientifically correct. Also, a limited sample cannot represent every resource on a platform. The evaluation must therefore explain how examples were selected and show the number of resources used for each result.

\clearpage
\section{Functional Requirements}

Functional requirements describe actions that the system must perform. Each requirement below can be checked with a command, database query, generated file, or test result. The requirements describe behaviour instead of individual function names, so they remain valid if the code is reorganized later.

\begin{table}[H]
\centering
\caption{Functional requirements of the extended pipeline.}
\label{tab:functional-requirements}
\small
\begin{tabularx}{\textwidth}{@{}p{0.9cm}p{9.5cm}p{2.2cm}@{}}
\toprule
ID & Requirement & Verification\tabularnewline
\midrule
FR1 & Detect GitHub, Codeberg, and Zenodo from a supported URL and reject unknown hosts. & Unit examples\tabularnewline
FR2 & Validate Git resources through Git and Zenodo resources through its record API before metadata is stored. & Connector test\tabularnewline
FR3 & Acquire Git repositories by clone or update and Zenodo records by file download and extraction. & Local files\tabularnewline
FR4 & Translate each supported API response into the same eight metadata fields without inventing missing values. & DB assertion\tabularnewline
FR5 & Support both one-resource input and platform-aware batch processing from SQLite. & End-to-end run\tabularnewline
FR6 & Register supplied notebook paths and discover \texttt{.ipynb} files automatically when paths are absent. & Notebook rows\tabularnewline
FR7 & Infer requirements, select Python, create an isolated environment, and execute Python notebooks. & Execution rows\tabularnewline
FR8 & Compare original and executed outputs and record status, metrics, errors, and elapsed time. & Result rows\tabularnewline
FR9 & Classify notebook purpose with rules and an AI model, store both attempts, and expose disagreement for review. & Classifier test\tabularnewline
FR10 & Export relational tables, materialise RDF with versioned mappings, and answer defined SPARQL questions. & KG test suite\tabularnewline
\bottomrule
\end{tabularx}
\end{table}

FR2 must happen before FR4. The pipeline checks that the source is valid before saving its metadata. Otherwise, information taken from an error response could be saved as if it were real metadata. The metadata is then fetched before the files are downloaded because the API already provides the source description and origin.

FR7 and FR8 reuse the original environment and execution steps. Each new connector only has to provide a local folder and notebook paths. It does not create a separate environment system. FR9 can run even when notebook execution fails because classification reads the notebook content without executing it. This allows failed notebooks to keep a category for later error analysis. FR10 requires the graph to be generated from the working database instead of edited by hand.

Together, these requirements describe one complete route from an input URL to results that can be queried. Downloading a file alone does not complete a connector: the resource must also be checked and described. Creating an N-Triples file alone does not complete the graph: its relationships and SPARQL test questions must also be correct.

\clearpage
\section{Non-Functional Requirements}

Non-functional requirements define the quality conditions under which the functional requirements must be fulfilled. They do not add new pipeline steps. Instead, they state how the pipeline must behave while collecting external resources, reading platform APIs, modifying database copies, executing notebooks, and publishing RDF results.

For this thesis, the most important quality goals are data safety, comparability, traceability, repeatability, robustness, security, and maintainability. These goals are necessary because a successful demonstration is not enough: the system must also protect the original data, explain how a result was produced, and remain understandable when new platforms or notebook samples are added.

\begin{table}[H]
\centering
\caption{Non-functional requirements of the extended pipeline.}
\label{tab:nonfunctional-requirements}
\footnotesize
\begin{tabularx}{\textwidth}{@{}>{\raggedright\arraybackslash}p{1.05cm}>{\raggedright\arraybackslash}p{2.45cm}>{\raggedright\arraybackslash}p{5.1cm}>{\raggedright\arraybackslash}X@{}}
\toprule
ID & Quality goal & Meaning in this thesis & Verification evidence\tabularnewline
\midrule
NFR1 & Data safety & Keep the source database unchanged and write new results only to the working copy. & Source database checksum or timestamp does not change.\tabularnewline
NFR2 & Comparability & Store platform results with shared fields, statuses, and metric definitions. & One query groups results by platform and status.\tabularnewline
NFR3 & Traceability & Link each result to platform, resource, run, execution, and classification method. & Foreign keys and RDF links connect the steps.\tabularnewline
NFR4 & Repeatability & Rebuild CSV and RDF outputs from the database, ontology, and mappings. & Generated files can be deleted and recreated.\tabularnewline
NFR5 & Robustness & Record expected failures without stopping the whole batch. & Failed resources keep status rows; later resources continue.\tabularnewline
NFR6 & Security & Treat URLs, archives, API text, tokens, and notebooks as untrusted. & Tokens are not logged; text is escaped; execution is isolated.\tabularnewline
NFR7 & Maintainability & Keep platform logic in connectors and notebook processing in shared stages. & Connector changes do not rewrite execution or RDF export.\tabularnewline
\bottomrule
\end{tabularx}
\end{table}

NFR1 protects the research baseline. The inherited GitHub database is input evidence, so the extension may read it but must not overwrite it. The first pipeline run creates a working copy in \texttt{data/output/db/}; all schema migrations, metadata rows, execution rows, and classification rows are then written to that copy. This makes experiments reversible because a new run can start again from the unchanged source database.

NFR2 and NFR3 make the results interpretable. Comparability does not mean that the platforms are treated as identical. It means that equivalent information, such as title, licence, run status, and output-comparison metrics, is stored in the same structure. At the same time, provenance keeps important differences visible: Zenodo remains an archived record, GitHub and Codeberg remain Git repositories, and every result can be followed from platform to resource, run, notebook execution, metric row, and classification activity.

NFR4 separates source evidence from generated artefacts. The working database, YARRRML/RML mappings, and ontology are the sources for semantic publication. CSV exports, N-Triples fragments, the combined graph, and SPARQL result files are generated outputs. They should therefore be rebuildable instead of manually edited. This requirement supports later verification because reviewers can inspect both the source rules and the generated graph.

NFR5 ensures that failures remain useful evidence. A batch may contain an unreachable repository, a missing Zenodo file, a package-installation failure, or a notebook execution error. These cases should be stored with a clear status and message so that error analysis can include them. They should not erase earlier results or prevent later resources in the batch from being processed.

NFR6 limits the risks created by external input. API tokens are read from environment variables and must not be written to logs or database fields. Text returned by APIs is escaped before database insertion. Downloaded archives are extracted only inside the intended resource directory. Notebook execution uses separate Python environments to avoid dependency conflicts. This isolation improves operational safety, but it is not a complete security sandbox; a public service would require stronger filesystem and network restrictions.

NFR7 keeps the design understandable. Platform connectors handle validation, metadata access, identifier normalization, and acquisition. After a local resource directory exists, notebook discovery, environment reconstruction, execution, comparison, classification, and knowledge-graph export follow shared code paths. This makes the thesis easier to evaluate because platform differences and platform-independent processing are described separately.

\section{Input and Output Contracts}

An input or output contract defines the exact information passed between parts of the system. In this pipeline, the contract is simple: the beginning of the pipeline may receive a resource in two different ways, but both ways must create the same six internal fields before notebook processing starts.

\noindent\textbf{Single-resource mode.} The input is one URL typed by a person. The URL points to a GitHub repository, a Codeberg repository, or a Zenodo record. The output is the same URL split into the six fields shown in Table~\ref{tab:io-contract}.

\noindent\textbf{Batch mode.} The input is one resource row already stored in the working database. No URL is typed by a person. The output is again the same six fields, rebuilt from the row's saved platform, identifier, and path values.

The important point is that the rest of the pipeline does not need to know which mode was used. It receives a URL, normalized identifier, platform, notebook paths, setup paths, and requirement paths. If some paths are unknown, their fields stay empty. This is common for Zenodo, because notebook paths are often only known after the record has been downloaded and extracted.

A GitHub or Codeberg notebook link needs one extra conversion. For example, a GitHub link ending in \texttt{/blob/main/analysis.ipynb} is a browser link, not a repository-relative path. The pipeline removes the display part, such as \texttt{/blob/main/} or \texttt{/src/branch/main/}, and keeps only the notebook path, for example \texttt{analysis.ipynb}. A Zenodo link is different: it identifies a record first, and the notebook paths are discovered later from the downloaded files.

\begin{table}[H]
\centering
\caption{Six internal fields produced by both input modes.}
\label{tab:io-contract}
\small
\begin{tabularx}{\textwidth}{@{}>{\raggedright\arraybackslash}p{3.2cm}>{\raggedright\arraybackslash}X>{\raggedright\arraybackslash}p{4.1cm}@{}}
\toprule
Field & Meaning & Example or empty case\tabularnewline
\midrule
URL & Original address used to validate and acquire the resource. & GitHub, Codeberg, or Zenodo URL\tabularnewline
Normalized identifier & Short platform-specific identifier used for database lookup. & \texttt{owner/notebooks} or \texttt{123456}\tabularnewline
Platform & Source platform detected from the URL or stored database row. & \texttt{github}, \texttt{codeberg}, or \texttt{zenodo}\tabularnewline
Notebook paths & Repository-relative notebook files selected for execution. & \texttt{analysis.ipynb}; empty means discover all notebooks.\tabularnewline
Setup paths & Optional setup files or scripts supplied with the resource. & \texttt{setup.py}; empty means no setup path was supplied.\tabularnewline
Requirement paths & Optional dependency files supplied with the resource. & \texttt{requirements.txt}; empty means infer dependencies from available evidence.\tabularnewline
\bottomrule
\end{tabularx}
\end{table}

These six fields are the handover from input handling to notebook processing. The URL tells the connector what to check and download. The normalized identifier gives the database a short lookup key. The platform prevents confusion between resources with similar names. The three path fields tell the pipeline which notebooks, setup files, or requirement files to use. If a path field is empty, the pipeline discovers the files itself when possible, or continues without that optional input.

After the six fields exist, notebook processing can start. The pipeline then stores its results in two places. The working database is the main evidence: it stores resources, metadata, runs, notebooks, executions, metrics, errors, and classifications. The file system stores supporting artefacts, such as logs, executed notebooks, and JSON comparison files. The knowledge-graph step reads the database, exports selected rows to CSV, and rebuilds the RDF graph from those CSV files and mapping rules. This means the database must be preserved, while the CSV and RDF files can be recreated from it.

\clearpage
\section{Data and Classification Requirements}

The database must separate a resource from the actions performed on it. The same repository or Zenodo record can be processed several times. One run can execute several notebooks, and one notebook can have several classification attempts. If all results were stored in one repository row, a new run would overwrite the old history.

\begin{table}[H]
\centering
\caption{Required relational entities and their purpose.}
\label{tab:data-entities}
\small
\begin{tabularx}{\textwidth}{@{}p{3.6cm}X@{}}
\toprule
Entity & Required meaning\tabularnewline
\midrule
\texttt{repositories} & Stable platform-specific resource identity and supplied path information\tabularnewline
\texttt{repository\_metadata} & One refreshable normalized description linked to the resource\tabularnewline
\texttt{notebooks} & Repository-relative notebook identity and language\tabularnewline
\texttt{repository\_runs} & One processing attempt with status, message, start, end, and duration\tabularnewline
\texttt{notebook\_executions} & Notebook-level execution, error, timing, and cell evidence\tabularnewline
\texttt{notebook\_}\linebreak\texttt{reproducibility\_}\linebreak\texttt{metrics} & Original-versus-executed comparison measurements\tabularnewline
\texttt{notebook\_}\linebreak\texttt{classifications} & Method-specific label, confidence, explanation, model/rule identity, and review state\tabularnewline
\bottomrule
\end{tabularx}
\end{table}

Normalized metadata has one current row for each resource. A later metadata fetch can update that row. Runs, executions, and classifications work differently: every new attempt receives a new row so that earlier results are not deleted. Foreign keys connect the tables, and a uniqueness rule prevents the same notebook path from being added twice to one repository.

The category list must be clear enough for people to use consistently while still covering common notebook purposes. The nine categories are \emph{data preparation}, \emph{data analysis}, \emph{visualization}, \emph{machine learning}, \emph{simulation}, \emph{tutorial}, \emph{software development}, \emph{mixed purpose}, and \emph{uncertain}. Every automated answer must select one of these labels and return a defined data structure, not only free text.

The rules and AI model should examine the same limited information: filename, notebook metadata, selected Markdown and code, imports, cell counts, and output types. Large outputs and possible secrets are not included unless necessary. The pipeline first stores the answer from each method. If the two answers agree, the result can be accepted automatically. Disagreement, an invalid answer, or low confidence is marked for human review. Manual evaluation labels are stored separately so they cannot be confused with predictions.

\clearpage
\section{Evaluation Requirements}

The evaluation must test both the software and the research questions. A demonstration with one repository shows that the route can work, but it is not enough evidence for a thesis comparison. The study also needs connector tests, database checks, classification evaluation, graph tests, and a documented experiment across all three platforms.

\begin{table}[H]
\centering
\caption{Evaluation matrix for the completed system.}
\label{tab:evaluation-matrix}
\small
\begin{tabularx}{\textwidth}{@{}p{3.1cm}p{4.3cm}X@{}}
\toprule
Area & Measures & Required evidence\tabularnewline
\midrule
Connector correctness & validation and acquisition success & Public examples and expected local files\tabularnewline
Metadata normalization & field accuracy and completeness & API-to-database checks for each platform\tabularnewline
Reproducibility route & environment, execution, and output rates & Structured run and notebook records\tabularnewline
Error analysis & counts by stage and category & Logs linked to classified failures\tabularnewline
Classification & accuracy, macro F1, per-class results, agreement & Human-labelled evaluation sample\tabularnewline
Knowledge graph & parseability, expected links, query answers & Triple counts and SPARQL competency tests\tabularnewline
End-to-end behaviour & repeatability and source-DB safety & Clean-run test and file checksums/schema checks\tabularnewline
\bottomrule
\end{tabularx}
\end{table}

The cross-platform experiment must explain how its examples were selected. It must also record the date, software versions, network conditions, and resource identifiers. For each platform, it should report how many resources were selected, validated, and downloaded; which metadata fields were present; how many notebooks were found; and how many environments, executions, and output comparisons succeeded. Percentages must always be shown together with the original counts, especially when the Codeberg and Zenodo samples are smaller than the GitHub data.

Classification must be checked against categories assigned manually without first showing the reviewer the automated answers. Overall accuracy is not enough because some categories may contain far more notebooks than others. The evaluation should also report precision, recall, macro F1, a confusion matrix, and examples of disagreement. The rules, AI model, and combined decision must be measured separately. This shows whether combining them actually improves the result.

Knowledge-graph evaluation must check more than whether the file opens. The ontology and mappings must be valid, CSV row counts must match the database queries, and the graph must contain the expected links between resources, runs, executions, results, and classifications. SPARQL questions test whether these links can be used correctly. A clean rerun must also leave the source database unchanged and recreate all documented outputs. The final claims are therefore supported by stored evidence, not only screenshots.
